\documentclass{article}
% {---
\usepackage[margin=1.0in]{geometry}

\setlength{\parindent}{0em} % default is 15 pt.
\setlength{\parskip}{5ex}
% \setlength{\headsep}{0pt}
% \setlength{\topskip}{0pt}
% \setlength{\topmargin}{0pt}
% \setlength{\topsep}{0pt}
% \setlength{\partopsep}{0pt}

\usepackage{enumitem}
\usepackage{tasks}
\usepackage{ragged2e}
\usepackage{graphicx}
\usepackage{xcolor}
\usepackage{multicol}
\usepackage{lipsum}
\usepackage{framed}
%\usepackage{tcolorbox}
\usepackage{ulem}  % \sout{text to cross out}
\usepackage{amsmath}
\usepackage{mathtools} % for "mathclap" in summation
\usepackage{mathptmx}  % \fontsize{12}{14}\selectfont ; times (including math)
\usepackage{mathpazo}  % palatino, including math

%\usepackage{fancyhdr}
%\pagestyle{headings}  % Put section name in top left, page number in top right
%\pagestyle{fancy}
%\fancyhf{}  % Clear all headers and footers (including default page number).

%\setlength{\headheight}{15pt}
%\renewcommand{\headrulewidth}{0pt} % remove the header rule
%\lhead{text} % Top left
%\rhead{text} % Top right
%\cfoot{text} % Top center
%\rfoot{\thepage} % Page numbers on bottom right corner
%\lfoot{\today} % current date on bottom left corner

\definecolor{mygray}{rgb}{0.52, 0.52, 0.51}
\definecolor{cadmiumorange}{rgb}{0.93, 0.53, 0.18}
\definecolor{darkcyan}{rgb}{0.0, 0.55, 0.55}
\definecolor{cadmiumgreen}{rgb}{0.0, 0.42, 0.24}
\definecolor{halayaube}{rgb}{0.4, 0.22, 0.33}
\definecolor{darkpowderblue}{rgb}{0.0, 0.2, 0.6}

%\newcommand{\mynotes}[1]{\textcolor{cadmiumorange}{#1}}
%\usepackage{tikz}\usetikzlibrary{backgrounds, shapes.misc}
%\newcommand*\circled[1]{\tikz[baseline=(char.base)]
%            \node[draw=black,shape=rounded rectangle,draw,inner sep=10pt] (char) {#1};}

\usepackage{titlesec}
%\setcounter{secnumdepth}{3}
%\usepackage{sectsty}
%\titleformat{<command>}
%   [<shape>]{<format>}{<label>}{<sep>}{<before-code>}[<after-code>]
\titleformat{\section}%
    [hang]
  {\Large\scshape} %\filcenter
  {\arabic{section}.\;}
  {10pt}     % sep
  {}        % before code
\titleformat{\subsection}%
  {\large\scshape} %\filcenter
  {\arabic{section}.\arabic{subsection}\;}
  {10pt}
  {}
\titleformat{\subsubsection}%
  {\scshape} %\filcenter
  {\arabic{section}.\arabic{subsection}.\arabic{subsubsection}\;}
  {10pt}
  {}

%\titlespacing*{⟨command⟩}{⟨left⟩}{⟨before-sep⟩}{⟨after-sep⟩}[⟨right-sep⟩]
%\titlespacing{\section}{}{}{}

% Section references/labels
%\renewcommand{\thesection}{}
%\renewcommand{\thesubsection}{\arabic{subsection}}
%\renewcommand{\thesubsubsection}{\arabic{subsubsection}}

% Lists
%\setlist[itemize]{noitemsep}
%\setlist[description]{%
%    itemsep=0ex,
%    align=right,
%    labelindent=,
%    labelwidth=,
%    labelsep=,
%    leftmargin=,
%}

% Found this in ~/Templates. No idea where it came from,
  % but looks interesting!
%\begin{lstlisting}
%
%\documentclass{article}
%\usepackage{pifont,xcolor}% http://ctan.org/pkg/{pifont,xcolor}
%\definecolor{myblue}{RGB}{0,29,119}
%\newcommand{\itemcolor}[1]{% Update list item colour
%                            \renewcommand{\makelabel}[1]{\color{#1}\hfil ##1}}
%\begin{document}
%
%\begin{dinglist}{110}
%\itemcolor{myblue}
%\item Sample text
%    \begin{dinglist}{228}
%    \itemcolor{green!70}
%    \item Sample text
%    \itemcolor{red!50}
%    \item Sample text
%        \begin{dinglist}{227}
%        \item Sample text
%        \end{dinglist}
%    \end{dinglist}
%\end{dinglist}
%
%\end{document}
%
%\end{lstlisting}


\setlist{noitemsep}
\renewcommand{\labelitemi}{$\vcenter{\hbox{\tiny$\bullet$}}$}
%\renewcommand{\labelitemi}{{\tiny$\bullet$}}
%   could just add this to \setlist[itemize]{label=...}
\definecolor{cadet}{rgb}{0.33, 0.41, 0.47}
%\renewcommand{\descriptionlabel}[1]{\ttfamily\textcolor{cadet}{#1}}
%\renewcommand{\descriptionlabel}{\ttfamily}
\setlist[description]{font={\ttfamily\color{halayaube}}}


% Verbatim
\usepackage{fancyvrb}  % framebox around verbatim text
\makeatletter
\renewcommand\verbatim@font{\normalfont\small\ttfamily}
\makeatother
% Change color of text inside \verb|text|
\usepackage{newverbs}
\newverbcommand{\cverb}{\color{red}}{}

\usepackage{listings}
\lstset{
  language=TeX,
  escapeinside={\%*}{*)},
  basicstyle=\small\linespread{1.0}\rmfamily,
  keywordstyle=\color{halayaube}\bfseries,
  commentstyle=\color{mygray},
%  frame=single,
  }

% Table of contents
\usepackage{setspace} % spacing between toc items
\usepackage[toc]{multitoc}
\renewcommand*{\multicolumntoc}{2}
%\setlength{\columnseprule}{0.5pt}

% Add dots to section level in table of contents
\makeatletter
\renewcommand*\l@section{\@dottedtocline{1}{1.5em}{2.3em}}
\makeatother

\usepackage{hyperref}
\hypersetup{
    colorlinks=true,
    urlcolor=blue,
    linkcolor=black,
}
%\urlstyle{same}

\title{LaTeX reference}
\author{Laurel Farris}
\date{\today}

% ---}


\begin{document}
\maketitle
\begin{itemize}
    \item \url{https://www.sharelatex.com/learn/Main_Page}
    \item \url{http://texdoc.net/texmf-dist/doc/latex/lshort-english/lshort.pdf}
    \item \url{http://texdoc.net/texmf-dist/doc/latex/titlesec/titlesec.pdf}
\end{itemize}

\setlength{\columnsep}{30pt}
\setcounter{tocdepth}{1}
\setlength{\parskip}{2pt}
%\begin{multicols*}{2}    % Start 2-columns
    %\addtocontents{toc}{\protect\setstretch{0.1}}
    \tableofcontents
%\end{multicols*}
\clearpage


\section{Units}
\begin{description}
    \item [px] pixels, depends on browser, use for electronic media
    \item [pt] points, use in print media
    \item [em] \emph{Horizontal} size, 1em is equal to the font size of
        the text.
    \item [ex] \emph{Vertical} size, 1ex is equal to the height of the
        letter `x' in the relevant font (usually).
\end{description}


\section{Line spacing}

\url{http://www.terminally-incoherent.com/blog/2007/09/19/latex-squeezing-the-vertical-white-space/}

\begin{lstlisting}
\usepackage[doublespacing]{setspace} % Double-spaced text
\renewcommand{\baselinestretch}{1.5} % 1.5 spacing (all text)
\usepackage{setspace}...\setstretch{1.5} % Apply to part of text?

\setlength{\parindent}{0pt} % 15 pt = default
\setlength{\parskip}{0pt}
\setlength{\headsep}{0pt}
\setlength{\topskip}{0pt}
\setlength{\topmargin}{0pt}
\setlength{\topsep}{0pt}
\setlength{\partopsep}{0pt}
\vspace{20 mm} % vertical blank space
    %\[small|med|big]skip
    %\[small|med|big]break
        % break same as skip, but also terminates paragraph and removes
        % preceeding v-space if bigger than would have been added by default
\hspace{20 mm} % horizontal blank space
\hfill % Pad with horizontal space to end of line
\hspace*{\fill}
\noindent
\,
\thinspace
\! % negative thin space
\: % medium space
\; % large space
\enspace
\quad
\qquad
\\ % Force line break
\mbox{text} % Prevent line break
\newline % Same as \\, but more vertical blank space?
\end{lstlisting}
Also see \S\ref{lists}.

\verb|\usepackage{ragged2e}|

\begin{minipage}{0.5\textwidth}
Commands:
\begin{lstlisting}
\flushleft
\enter
\flushright
\end{lstlisting}
\end{minipage}
\begin{minipage}{0.5\textwidth}
Environments:
\begin{lstlisting}
\begin{flushright}...\end{flushright}
\begin{flushleft}...\end{flushleft}
\begin{center} ... \end{center}
\begin{justify} ... \end{justify}
\end{lstlisting}
\end{minipage}

NOTE: Vertical space commands won't work if they're part of a horizontal line.
E.g. \verb|\vfill| and \verb|\vspace| must be preceeded by a line break,
and there needs
to be physical content on both sides of the space.
Use \verb|\null| to create `fake content', e.g.
\verb| \newpage \null \vfill ... |

\subsection{Page breaks}

\begin{lstlisting}

\newpage{
    % Technically new COLUMN:
    %   For single-column documents (probably anything that isn't an apj article),
    %   a new column IS a new page.
    %   If using multicol environment, starts new column, possibly on same page.
    %   EXCEPTION: if using multicol PACKAGE, starts new page every time.
    %   (NOTE: use \columnbreak for new column with multicol package)

\clearpage
  % Like \newpage, but also restricts floats.
  % Useful for placing figures where you want them
  %   (had to put this before and after each deluxetable in aastex).
  % "Forces pending floats from the stack to be typeset starting on the page
  %   that follows the page break. Flushes out all floats that have been deferred."

\begin{samepage}... \end{samepage} % PREVENT page break.

\end{lstlisting}


\subsection{Margins}
\url{https://tex.stackexchange.com/questions/42371/margins-and-text-position}

\begin{lstlisting}
\usepackage{fullpage}

\addtolength{\oddsidemargin}{-0.875in}
\addtolength{\evensidemargin}{-0.875in}
\addtolength{\textwidth}{1.75in}
\addtolength{\topmargin}{-0.875in}
\addtolength{\textheight}{1.75in}

% To increase the inner margins, i.e. for 3-prong folders:
\addtolength{\oddsidemargin}{0.5in}
\addtolength{\evensidemargin}{-0.5in}

\usepackage[options]{geometry}
% or set options separately:
\usepackage{geometry}
\geometry{options}

\geometry{
    margin=1in,
    left=1in, right=1in, top=1in, bottom=1in,
    textwidth=6.0in,
    paperwidth=140mm, paperheight=105mm,
    hmargin=1cm, vmargin=1cm,
}

% Set bindingoffset option to shift inner margins
\usepackage[a4paper,width=150mm,top=25mm,bottom=25mm,bindingoffset=6mm]{geometry}



% Change geometry within document
\newgeometry{left=3cm, bottom=0.1cm}
...
\restoregeometry

% Change geometry using changepage package and "adjustwidth" environment.
\usepackage{changepage}
\begin{document}
    ...
\begin{adjustwidth}{<left>}{<right>}
    ...
\end{adjustwidth}
% NOTE: <left> and <right> values are ADDED to existing margins.

% Custom environment to change margins in only a portion of text;
% it will indent the left and right margins by the values given.
\newenvironment{changemargin}[2]{%
 \begin{list}{}{%
  \setlength{\topsep}{0pt}%
  \setlength{\leftmargin}{#1}%
  \setlength{\rightmargin}{#2}%
  \setlength{\listparindent}{\parindent}%
  \setlength{\itemindent}{\parindent}%
  \setlength{\parsep}{\parskip}%
 }%
\item[]}{\end{list}}
\end{lstlisting}



\clearpage
\section{Headers and footers}\label{headfoot}
\begin{minipage}{0.5\textwidth}
\begin{lstlisting}
\usepackage{fancyhdr}

% Automatically generates a header
%   with section name
\pagestyle{fancy}

\setlength{\headheight}{15pt}
\lhead{text} % Top left
\rhead{text} % Top right
\chead{text} % Top center
\lfoot{text} % Bottom left
\rfoot{text} % Bottom right
\cfoot{text} % Bottom center
    % text can be \today for current date or \thepage for page #s.

%% ---
% Headers & Footers (from Dissertation, 3/13/2020)
%
%\usepackage{fancyhdr}
%\setlength{\headheight}{15.0pt}
%\pagestyle{fancy}
%\fancyhf{}
%\renewcommand{\headrulewidth}{0.5pt} % remove the header rule
%\fancyhead[R]{\footnotesize\today}
%\renewcommand*{\headrule}{\color{black!50!white}}
%
% Change format of date (sorta... still year month day, just changes separator.
%\usepackage{yyyymmdd]{datetime}
%\renewcommand{\dateseparator}{−−}
%\renewcommand{\dateseparator}{−}
%\renewcommand{\dateseparator}{/}
%
%\cfoot{\today}
%% ---


\end{lstlisting}
\end{minipage}%
\begin{minipage}{0.5\textwidth}
    \begin{center}
        The \verb|\headheight| option sets the amount of space between the
        header and the top edge of the paper. Value has to be greater than
        13.6, otherwise will get an error message. Document still
        compiles, but better safe than sorry. Setting the left, center, and/or
        right headers overwrites the one generated automatically.
    \end{center}
\end{minipage}


\subsection{Page numbers}
\begin{lstlisting}
\pagenumbering{gobble}
\pagestyle{empty}
\fancyhf{}  % Clear all headers and footers (including default page number).
\renewcommand{\headrulewidth}{0pt} % remove the header rule
\rfoot{\thepage}
\lfoot{\thepage}
\end{lstlisting}

\subsection{Footnotes}
\begin{minipage}[t]{0.5\textwidth}
\begin{lstlisting}
\usepackage[stable,symbol]{footmisc}
% stable: footnotes in section titles
% symbol: symbols instead of numbers

% Re-start markers after each page:
\usepackage{perpage}
\MakePerPage{footnote}
\end{lstlisting}
\end{minipage}%
\begin{minipage}[t]{0.5\textwidth}
    Example:
\begin{lstlisting}
Here is some relevant information\footnote{
    See Guy et al. for additional information.}.
\end{lstlisting}

Here is some relevant information\footnote{
    See Guy et al. for additional information.}.
\end{minipage}

\begin{lstlisting}
\renewcommand{\footnoterule}{%
  \kern -3pt
  \hrule width \textwidth height 1pt
  \kern 2pt
}
\renewcommand\footnoterule{\rule{\linewidth}{5pt}} % Another option
\end{lstlisting}

\begin{minipage}{0.5\textwidth}
\begin{lstlisting}
% Current date (month, dd, yyyy)
%   in left|right|center footnote
\fancyfoot[L|R|C]{\tiny\itshape\today}
\end{lstlisting}
\end{minipage}%
\begin{minipage}{0.5\textwidth}
\begin{lstlisting}
% Change format of date
\usepackage{yyyymmdd]{datetime}
\renewcommand{\dateseparator}{--}
\renewcommand{\dateseparator}{-}
\renewcommand{\dateseparator}{/}
\end{lstlisting}
\end{minipage}

\newpage
\section{Fonts}
\begin{itemize}
    \item \href{https://www.tug.org/pracjourn/2006-1/schmidt/schmidt.pdf}{
            The \LaTeX{} font catalogue}
    \item \href{https://en.wikibooks.org/wiki/LaTeX/Fonts}{
            Wikibooks -- \LaTeX{} fonts}
\end{itemize}

Font that applies to entire document:

\begin{lstlisting}
\usepackage{lmodern}
\renewcommand\familydefault{\sfdefault} % base font of the document
\renewcommand*\familydefault{\sfdefault} % Difference from above??
\usepackage[T1]{fontenc}
\end{lstlisting}

NOTE: bfseries overrides scshape when using lmodern package.

\subsection{Font size}

\begin{lstlisting}
\documentclass[12pt]{article}
\documentclass[11pt]{article}
\documentclass[10pt]{article} % default
\end{lstlisting}

\begin{lstlisting}
\fontsize{<size>}{<baselineskip>}\selectfont
\end{lstlisting}

The baseline-skip should be set to roughly 1.2x the font size.

\begin{minipage}[t]{0.3\textwidth}
\begin{lstlisting}
\Huge
\huge
\Large
\large
\normalsize
\small
\footnotesize
\scriptsize
\tiny
\end{lstlisting}
\end{minipage}%
\begin{minipage}[t]{0.7\textwidth}
    Example:
\begin{lstlisting}
{\Large I want this text to be big.}
\end{lstlisting}
    \vspace{-2ex}{\Large I want this text to be big.}\\
(enclosing entire thing in \verb|{}|s keeps from having to use
\verb|\normalsize| at the end).

Can also use environments:
\begin{lstlisting}
\begin{Large}
I want this text to be big.
\end{Large}
\end{lstlisting}
\end{minipage}

\subsection{Font style}
\subsubsection{Modal}

\begin{minipage}[t]{0.5\textwidth}
\begin{lstlisting}
\mdseries
\bfseries
\upshape
\itshape
\scshape
\slshape
\rmfamily
\sffamily
\ttfamily
\end{lstlisting}
\end{minipage}%
\begin{minipage}[t]{0.5\textwidth}
    These don't read text as an argument, and can somehow be
    used in the verbatim environment?
\end{minipage}

\subsubsection{Textblock}
\begin{minipage}[t]{0.5\textwidth}
\begin{lstlisting}
\textbf{bold}
\textit{italics}
\texttt{computer style}
\textsf{sans serif}
\textsl{slanted}
\textsc{small caps}
\emph{emphasis}
\underline{underlined}
\sout{strikethrough}
\end{lstlisting}
\end{minipage}%
\begin{minipage}[t]{0.5\textwidth}
    \begin{list}{}
        \item \textbf{bold}
        \item \textit{italics}
        \item \texttt{computer style}
        \item \textsf{sans serif}
        \item \textsl{slanted}
        \item \textsc{small caps}
        \item \emph{emphasis}
        \item \underline{underlined}
        \item \sout{strikethrough}
    \end{list}
\end{minipage}

Note: the \verb|ulem| package is required to cross out text using
\verb|sout|, and also causes \verb|\emph| to push emphasis on
text by underlining it, rather than putting it in italics
(ulem = \underline{u}nder\underline{l}ine \underline{em}phasis).

\clearpage
\section{Sections}
\begin{minipage}{\textwidth}
\href{https://www.sharelatex.com/learn/Sections_and_chapters#Numbered_and_unnumbered_sections}
    {Numbering}
\end{minipage}

\begin{lstlisting}
    % Set depth of sections to number
    \setcounter{secnumdepth}{0} % Don't number sections
    \setcounter{secnumdepth}{1} % Number sections
    \setcounter{secnumdepth}{2} % Number sections and subsections ...
\end{lstlisting}


\begin{minipage}[t]{0.5\textwidth}
\begin{lstlisting}
    \section{Title of first section}
    \subsection{Title of subsection}
    \subsubsection{Title of subsubsection}
    \paragraph{Paragraph heading}
    \subparagraph{Subparagraph heading}
\end{lstlisting}
\end{minipage}%
\begin{minipage}[t]{0.5\textwidth}
Paragraphs are not numbered or followed by a line break.
Note that \verb|\paragraph{}| and \verb|\par| are not the same thing.
\verb|\par| does the same thing as a blank line in the text file (starts a new paragraph).
\end{minipage}


\subsection{Customize section headings}
NOTE: This is for the headings only, not the following text.
Can't change vertical space between heading and text, or text color, since
starting a section isn't the same as starting an environment.

NOTE: if using aastex, the title, abstract, and keywords will not appear in
pdf if using the titlesec package.
\begin{lstlisting}
\usepackage{titlesec}
\titleformat{<command>}
    [<shape>]{<format>}{<label>}{<sep>}{<before-code>}[<after-code>]
\end{lstlisting}

\begin{minipage}{0.5\textwidth}
Shape:
\begin{lstlisting}
    hang % Default
    rightmargin, leftmargin
      % Titles are in the margins,
      %   rather than body of page.
\end{lstlisting}
\end{minipage}%
\begin{minipage}{0.5\textwidth}
Labels:
\begin{lstlisting}
    \arabic   (1, 2, 3, ...)
    \alph     (a, b, c, ...)
    \Alph     (A, B, C, ...)
    \roman    (i, ii, iii, ...)
    \Roman    (I, II, III, ...)
    \fnsymbol (∗, †, ‡, §, ¶, ...)
\end{lstlisting}
\end{minipage}

Can also simply tweak the existing title format:
\begin{lstlisting}
    \usepackage{titlesec}
    \titleformat*{\section}{\LARGE\bfseries}
    \titleformat*{\subsection}{\Large\bfseries}
    \titleformat*{\subsubsection}{\large\bfseries}

    \renewcommand{\thesection}{Text \arabic{section}}   % Text in front of label
    \renewcommand{\thesection}{\Roman{section}}         % Roman numerals
\end{lstlisting}

\verb|\par| can't be used to start a new line in \verb|titleformat|, but
\verb|\\| and
\verb|\newline| can.

\clearpage
Examples:
\begin{lstlisting}
\titleformat{\section}%
  {\fontsize{16}{18}\selectfont\bfseries\color{myblue}}
  {\fontsize{46}{50}\selectfont\color{mypur}\arabic{section}\color{black}$\vert$}
  {0em}{}
\titleformat{\subsection}%
  {\fontsize{14}{16}\selectfont\bfseries\color{mypur}}
  {\color{myblue}\circled{\arabic{section}.\arabic{subsection}}}
  {0.5em}{}
  [\vspace{-2.5pt}{\color{mygray}\titlerule[5pt]}]
  %[\vspace{-20pt}\colorbox{mygray}{
  %  \begin{minipage}{\textwidth}
  %    \vspace*{2pt} % Space before
  %    \hfill
  %    \vspace*{2pt} % Space after
  %  \end{minipage}}]
\titleformat{\subsubsection}%
  {\fontsize{13}{14}\selectfont\bfseries\color{mypur}}
  {\color{myblue}\arabic{section}.\arabic{subsection}.\arabic{subsubsection}}
  {1em}{}
  [\vspace{-2.5pt}{\color{mygray}\titlerule[3pt]}]
\newcommand\colonafter[1]{#1:}
\titleformat{\paragraph}% Add colon to end of all paragraph headings:
  {\fontsize{12}{13}\selectfont\bfseries\color{myblue}}
  {}
  {0.5em}
  {\colonafter}
\end{lstlisting}

\subsection{Space around section titles}
\begin{minipage}[t]{\textwidth}
\begin{lstlisting}
\usepackage{titlesec}
\titlespacing*{<command>}{⟨left>}{<before-sep>}{<after-sep>}[<right-sep>]
\titlespacing*{\section}{-0.50in}{0pt}{0pt}[]
\titlespacing*{\subsection}{-0.25in}{0pt}{0pt}[]
\titlespacing*{\subsubsection}{-0.25in}{0pt}{0pt}[]
\titlespacing*{\paragraph}{-0.25in}{0pt}{0pt}[]
\titlespacing*{\subparagraph}{-0.25in}{0pt}{0pt}[]
\end{lstlisting}
\end{minipage}
\begin{minipage}[t]{\textwidth}
    The asterisk removes paragraph indentation following the section title
    (doesn't do anything if there is no indentation anyway). It also appears to
    allow you to set only a few options in titleformat without creating empty
    braces for every single argument (adds to existing title format, rather
    than starting from scratch\ldots?)

    This is where you adjust vertical space. Don't try to do this in titleformat!
\end{minipage}

\subsection{Cross-referencing in text}
\begin{lstlisting}
See section \S\ref{data} for the data description.
...
\subsection{The Data}\label{data}
...
\end{lstlisting}
May need to run \texttt{pdflatex} twice for this to take effect.
Obviously won't have anything to refer to if the sections aren't numbered.
There should be a way to do this though, like by linking the section name.

\clearpage
\section{Table of contents}\label{contents}%

\url{http://texblog.org/2011/09/09/10-ways-to-customize-tocloflot/}\\
\url{http://tex.stackexchange.com/questions/37940/table-of-contents-with-roman-arabic-and-no-page-numbers}

\begin{lstlisting}
% n = number of levels deep to go.
\setcounter{tocdepth}{n}

\usepackage[toc]{multitoc}

% Multi-column toc (default: n=2)
\renewcommand*{\multicolumntoc}{n}

% Horizontal separation between toc columns
%  (print out as \the\columnsep)
\setlength{\columnsep}{30pt}

% Width of line between toc columns
\setlength{\columnseprule}{0.5pt}
\end{lstlisting}

Notes:
\begin{itemize}
    \item You will have to run \texttt{pdflatex} twice.
    \item It appears that creating a toc puts headers on all pages, which may
        not be desired. See \S{}\ref{headfoot} for getting rid of them.
    \item Need hyperref package in order to have clickable links in toc.
\end{itemize}

Vertical space between toc items:
\begin{lstlisting}
\usepackage{setspace}
...
\addtocontents{toc}{\protect\setstretch{n}}
    % ``protect'' has something to do with ``fragile'' things.

\renewcommand{\baselinestretch}{<n>}\normalsize
    % n=2 for double spacing, 1 for single, 0.75 for compress.

\setlength{\parskip}{0pt}  % Another way to change spacing
\tableofcontents
\setlength{\parskip}{10pt}  % Spacing for remainder of document
\end{lstlisting}

Some sections, like those with `*' won't be included. To add them:
\begin{lstlisting}
% Syntax:
\addcontentsline{type}{section_level}{entry}
% Example:
\addcontentsline{toc}{section}{Preface}
\end{lstlisting}


Include figures and tables in table of contents:
\begin{lstlisting}
\listoffigures
\listoftables
\setcounter{lofdepth}{2} % lof = list of figures
\end{lstlisting}
NOTE: figure and table environments need to be used.

Change heading in toc for each list:
\begin{lstlisting}
\renewcommand\contentsname{}
\renewcommand\listfigurename{}
\renewcommand\listtablename{}
\end{lstlisting}

Fill space between section and page number with dots.
By default this is only done for subsections and below.
\begin{lstlisting}
\makeatletter
\renewcommand*\l@section{\@dottedtocline{1}{1.5em}{2.3em}}
\makeatother
...
\begin{document}
\end{lstlisting}

minitoc?!?

\clearpage
\section{Margin notes}
Add notes to margins with \verb|\marginpar{text in margins}|,
or for more flexibility, use the \verb|marginnote| package:
\begin{lstlisting}
\usepackage{marginnote}
\usepackage{showframe,marginnote}  % box around margins
\setlength{\marginparwidth}{1in}
\renewcommand*{\marginfont}{\color{red}\sffamily}
\renewcommand*{\raggedleftmarginnote}{} % ?
\renewcommand*{\raggedrightmarginnote}{\centering} % ??
\renewcommand*{marginnotevadjust}{40pt} % Vertically shift margin notes (?)

\begin{document}
\marginnote{<right>} % aligned left (ragged right)
\marginpar{<right>} % aligned left (ragged right)
\reversemarginpar  % Switch to left side margins
\marginnote{<left>} % aligned right (ragged left)
\marginpar{<left>}  % aligned left (ragged right)
\normalmarginpar  % switch back

% Maybe don't need \reversemarginpar ?!?
\marginnote[<left>]{<right>}[vertical drop in cm]

% Example:
\hspace{0pt} % Put note next to section title text, rather than above it,
% which is done by default for some stupid reason
\reversemarginpar
\marginnote{}
\end{lstlisting}


Use curly braces to mark text that margin notes are referring to:

\begin{lstlisting}
\usepackage{schemata}
\schema[open|close]{body text}{margin text}
\end{lstlisting}
Sized to enclose whatever text it's opening toward (body or margin)
although same amount of space will be allotted for the other\ldots

Put a vertical bar between document text and notes in margin:
\begin{lstlisting}
\usepackage[color,leftbars]{changebar} % On right side by default
\cbcolor{red}
\setlength\changebarsep{10pt}
\setlength\changebarwidth{1pt}
...
\cbstart
...
\cbend
\end{lstlisting}

\clearpage
\section{Comments}

\href{http://mirrors.ibiblio.org/CTAN/macros/latex/contrib/comment/comment.pdf}{
    Comment package}

Section 3.1 explains the \texttt{comment.cut} file created when running
pdflatex\ldots should contain the text commented in the main \texttt{.tex} file.

\begin{lstlisting}

\usepackage{comment}
\includecomment{mycomment} % Include content in output pdf
\excludecomment{mycomment} % Exclude content, i.e. treat as real comment.
\specialcomment{mycomment}{
    \begingroup\sffamily\small\color{red!50!black}}{
    \endgroup}

\begin{document}
...
\begin{comment}
    Text to be commented\ldots
\end{comment}
...
\end{document}
\end{lstlisting}






\clearpage
\section{Lists}\label{lists}
%\begin{itemize}
%    \item \url{ftp://ftp.nsu.ru/mirrors/ftp.dante.de/tex-archive/macros/latex/contrib/enumitem/enumitem.pdf}
%    \item \url{https://www.ntg.nl/maps/11/33.pdf}
%    \item \url{https://www.sharelatex.com/learn/Lists#Reference_guide}
%    \item \url{http://ctan.mirrors.hoobly.com/macros/latex/contrib/enumitem/enumitem.pdf}
%    \item \url{http://www.troubleshooters.com/linux/lyx/ownlists.htm}
%\end{itemize}

\begin{lstlisting}
\usepackage{enumitem}

% Apply customizations globally:
\setlist*[<typeoflist>,<n>]{<option=value>} % Add settings to current
\setlist[<typeoflist>,<n>]{<option=value>} % Start from scratch

% Apply customizations to one indivdual list:
\begin{<typeoflist>}[options] ... \end{<typeoflist>}
\end{lstlisting}

%\setlist*[description]{font={\color{blue}}}
\begin{description}[
            labelwidth=23mm, labelsep=2mm, leftmargin=25mm,
            itemsep=2mm, align=right,]
    \item [typeoflist] = \verb|itemize|, \verb|enumerate|, \verb|description|,
        or \verb|list|. If not specified, options will be applied to all types.
    \item [n] = nested level (=1 for top level). If not specified, options
        will be applied to all levels.
\end{description}

Available \verb|options| are as follows:
%\setlist*[description]{font={\color{red}}}
\begin{description}[
            labelwidth=23mm, labelsep=2mm, leftmargin=25mm,
            itemsep=2mm, align=right,]
    \item [font] e.g. \verb|font={\bfseries\sffamily\color{red}}|
    \item [labelindent] (horizontal) Appears to be the width between edge of text
        and left side of label. Default must be a negative number, since
        setting this to 0in aligns the labels with the text.
    \item [labelwidth] (horizontal)
        Width allotted to the label. This should be equal to or greater than
        the longest \emph{expected} label. Good for lining up text when labels
        are left-aligned. This will override labelindent if order is switched!
    \item [labelsep] (horizontal)
        Space between label (right edge of labelwidth) and text
        (left margin of the item body).
        NOTE: The labelsep is enforced by shoving the label left rather
        than moving the body left margin right\ldots may need
        to add a corresponding amount to \texttt{leftmargin}.
    \item [leftmargin] (horizontal) Indentation of body text below label.
        Distance from the left edge of the label to the left margin of the item
        body. Remember, environments can nest. Defaults to 0. Eventually maxes
        out and returns to default value.
    \item [rightmargin] Indentation for right margin of body text.
    \item [listparindent] (horizontal) Indentation of the first line of each
        paragraph (except the first) in an item. Helpful for decreasing
        interparagraph spacing while cuing beginning of new paragraphs.
    \item [itemindent] (horizontal) ``Horrid adjustment''\ldots
        apparently this one should be avoided.
    \item [headsep] (vertical)
    \item [itemsep] (vertical) Extra inter-item spacing added to
        \texttt{parsep}. Set option to \textcolor{red}{\texttt{noitemsep}} for
        no additional space between items (also helps to even out the amount of
        space between items in nested lists).
    \item [topskip] (vertical)
    \item [topmargin] (vertical)
    \item [topsep] (vertical) Extra space added to \texttt{parskip} before the first
        AND after the last item\ldots bit of a misnomer.
        Also may not be just for lists\ldots
    \item [parskip] (vertical) Space between paragraphs outside of a list, and part of
        the space between a non-list paragraph and a list item.
        \textcolor{red}{This is NOT a list property; it can be set
        globally for entire document (see \SS{} \texttt{ref\{\}}).}
    \item [parsep] (vertical) Paragraph separation within a single item.
    \item [partopsep] (vertical) This is added to the top and/or bottom of the list
        if and only if there's a blank line above or below the first
        or last item. \emph{Leave this alone unless blank lines become a
        problem}. Also may not be just for lists\ldots
\end{description}

%NOTE: without \texttt{enumitem} package, alternative to \verb|noitemsep| is
%the \verb|mdwlist| package - add asterisks to list environments.
%\begin{lstlisting}
%\usepackage{mdwlist}
%...
%\begin{itemize*} ... \end{itemize*}
%\end{lstlisting}

\subsection{itemize}
Change label (bullets by default) size/style:
\begin{lstlisting}
\renewcommand{\labelitemi}{$\vcenter{\hbox{\tiny$\bullet$}}$}
\renewcommand{\labelitemi}{{\tiny$\bullet$}} % Same as above?
\begin{itemize}[label={}]  % No bullet at all

\end{lstlisting}

\subsection{enumerate}
Change label (numbers by default):
\begin{lstlisting}
\setlist[enumerate,1]{label=...}
% Or: \setenumerate[0]{label=...}   ?

% Labels:
    {(\arabic*)}  % Numbers
    {(\roman*)}   % Roman numerals
    {(\Alph*)}    % Letters --> (A), (B), ...
    {(\alph*)}    % Letters --> (a), (b), ...
% NOTE: the asterisk connects the physical level of the list
% (e.g. gives label `b' to second item when using \alph).

% 1.1, 1.2  -->  1.2.1, 1.2.2, etc.
\setlist[enumerate,1]{label={\arabic{section}.\arabic*}}
\setlist[enumerate,2]{label={\arabic{section}.\arabic{enumi}.\arabic*}}
\end{lstlisting}

Or use the \texttt{enumerate} package:
\begin{lstlisting}
\usepackage{enumerate}
\begin{document}
\begin{enumerate}[label*=\arabic*.]
\begin{enumerate}[I]
\begin{enumerate}[I.]
\begin{enumerate}[(a)]
\end{lstlisting}

\begin{lstlisting}
% Change value of first label:
\begin{enumerate}
    \setcounter{enumi}{<n>}
    \item ...
    \item ...
\end{enumerate}
% n = index
%   NOTE: array of possible starting numbers are indexed starting at 0!
%     So to begin enumerated list with 5, set n = 4.
% \setcounter must be inside enumerate environment.
\end{lstlisting}

\clearpage
\subsection{description}
To customize the description labels (the items inside the brackets),
in the preamble:
\begin{lstlisting}
\renewcommand{\descriptionlabel}[1]{%
    \hspace{\labelsep}
    \ttfamily
    \textcolor{red}{#1}
}
\end{lstlisting}
This puts the labels in typewriter font in a different color.
By default, description labels start a distance equal to
\texttt{hspace} to the \emph{left} of the text, so adding that line
causes them to line up with the left edge of the text instead.

\begin{lstlisting}
align=right
leftmargin=*        % Align with main text... value affects text after the first line
font=\normalfont    % Not bold, which is the default
style=nextline      % Description starts on the next line
style=multiline     % ???
\end{lstlisting}

If you want the label itself to include brackets, put braces around them:
\begin{lstlisting}
\item [ label text {[label text in brackets]} ] description text
\end{lstlisting}


\subsection{list}
\begin{lstlisting}
\begin{list}{default_label}{decls}
    default label: Text to be used as a label (leave blank if none desired)
    decls: geometrical parameters
\end{lstlisting}

\subsection{tasks}
\begin{minipage}{0.5\textwidth}
    \begin{lstlisting}
    \usepackage{tasks}
    ...
    \begin{tasks}(4)
        \task one
        \tast two
    \end{tasks}
    \end{lstlisting}
\end{minipage}%
\begin{minipage}{0.5\textwidth}
    \begin{tasks}(4)
        \task one
        \task two
    \end{tasks}
\end{minipage}

These will be listed horizontally, rather than vertically.
The number in parentheses affects the horizontal space between each task;
the higher the number, the \emph{smaller} the spacing. Maybe refers to how
many tasks should fit on one line?


\clearpage
\section{Colors}\label{color}

\begin{lstlisting}
\usepackage{color}
\usepackage{xcolor}
\end{lstlisting}

The \texttt{color} package is needed to use pre-defined colors (white, black,
red, green, blue, cyan, magenta, yellow). The \texttt{xcolors} package is
needed to define new colors (see \S{}~\ref{definecolors}), color the background,
and to use color mixtures.

\begin{lstlisting}
% Color background
\pagecolor{yellow!30}

% Applied to a small bit of text.
\textcolor{green}{I want this text to be green.}
\textcolor[rgb]{0,1,0}{text}  % green

% applied to all following until color is changed again
%   or inside environment containing the statement
\color[rgb]{1,0,0} % red

\colorlet{<new color name>}{<old color name>}

%  30% blue, 70% green
\color{blue!30!green}

%  0.20(0,0,1) +
%  (1-0.20)(1,0,0) +
%  0.3(1,0,0) +
%  (1-0.30)(0,1,0)
\color{blue!20!red!30!green}

\end{lstlisting}

\subsection{Define your own colors}\label{definecolors}
\url{http://latexcolor.com}
\begin{lstlisting}
\usepackage[usenames, dvipsnames]{color}
 \definecolor{color}{HTML}{AF00D7} % HTML must be in caps!
 \definecolor{mypink1}{rgb}{0.858, 0.188, 0.478}
 \definecolor{mypink2}{RGB}{219, 48, 122}
 \definecolor{mypink3}{cmyk}{0, 0.7808, 0.4429, 0.1412}
 \definecolor{mygray}{gray}{0.6}
\textcolor{mygray}{text I want to be gray}.
\end{lstlisting}

\clearpage

\section{Hyperlinks}
In preamble:
\begin{lstlisting}
% break long links over multiple lines.
\usepackage[breaklinks=true]{hyperref}
\hypersetup{
    hidelinks,         % ?
    linktoc=all,       % ?
    linktocpage=false, % only page #s are clickable links in toc
    hyperindex=true,   % Makes the page numbers of index entries into hyperlinks
    breaklinks=false,  % Allows links to be broken into multiple lines.
    frenchlinks=false, % Use small caps instead of colours for links
    colorlinks=false,  % Allow color for link and anchor text...
    urlcolor=magenta   %   linked URLs
    linkcolor=red      %   normal internal links
    citecolor=green,   %   bibliographical citations
    anchorcolor=black, %   anchor (target) text
    filecolor=cyan,    %   links that open local files
}


% Make \href{} link text bold
\let\oldhref\href
\renewcommand{\href}[2]{\oldhref{#1}{\bfseries#2}}

% Change style of urls (same as body text, roman, sans serif, bold, etc.)
%  NOTE: url text is not the same as hyperlink text in \href{}.
    \urlstyle{same}
    \urlstyle{rm}
    \urlstyle{sf}
    \bfseries\urlstyle{same}
    \itshape\urlstyle{same}
    \mdseries\itshape\urlstyle{same}

\end{lstlisting}

\href{https://www.overleaf.com/learn/latex/Hyperlinks#Reference_guide}{
    hypersetup options}


NOTE: `draft' option (\verb|\documentclass[draft]{...}|) removes toc links.

Insert hyperlink in text:
\begin{lstlisting}
\url{http://google.com}
\href{http://google.com}{link text}
\end{lstlisting}

Insert link to local file:
\begin{lstlisting}
\href[]{run:/path/to/paper.pdf}{text}
\end{lstlisting}

Link text to other text in document:
\begin{lstlisting}
    \hypertarget{word1label}{\hyperlink{word2label}{Word1}}
    \hypertarget{word2label}{\hyperlink{word1label}{Word2}}
\end{lstlisting}

\clearpage
\section{Textboxes}
Several options:
\begin{lstlisting}
% Simple frame
\usepackage{framed}
\begin{framed}...\end{framed}

% colorbox (does not support line breaks)
\colorbox{hl}{\parbox{0.9\textwidth}Text in box}

% fcolorbox
\fboxrule=<value><unit> % Adjust thickness of border
\fboxsep=<value><unit> % Adjust thickness of padding
\fcolorbox{<border-color>}{<background-color>}{<contents>}

% Use \textrm{t}, \textrm{c}, \textrm{b} options
% to align the base line of the top, center, and bottom row
% with the surrounding baseline.

% Example:
\fboxrule=1pt and \fboxsep=5pt
Text...
\medskip % ?
\noindent\fcolorbox{red}{yellow}{%
   \minipage[t]{\dimexpr0.48\linewidth-2\fboxsep-2\fboxrule\relax}
      Text...
   \endminipage}\hfill
   \fcolorbox{red}{yellow}{%
   \minipage[t]{\dimexpr0.48\linewidth-2\fboxsep-2\fboxrule\relax}
      Text...
   \endminipage}
\medskip
Text...
\end{lstlisting}

\clearpage
\subsection{tcolorbox}
\url{http://mirrors.ibiblio.org/CTAN/macros/latex/contrib/tcolorbox/tcolorbox.pdf}

\begin{lstlisting}
\usepackage{tcolorbox}
\tcbset{
    nobeforeafter, % place boxes side by side
    colback=color, % background color of entire box
    colbacktitle=color,  % background color of title
    colframe=color, % color of frame around box
    coltitle=color, % color of title text
    detach title, % ?
    before={\vspace{1in}}, % vertical space between boxes
    before upper={\hfill\tcbtitle\par\vspace{0.5in}}, % right aligned title, upper part of box (if it's divided into upper/lower sections)
    fonttitle={\Large\bfseries,nobeforeafter,center title},
    fontupper|fontupper=\fontsize{14pt}{16pt}\selectfont,
    width=4cm,
    height=8cm,
    boxrule=3mm, % width of all four sides
    toprule=3mm,
    bottomrule=3mm,
    leftrule=3mm,
    rightrule=3mm,
    arc=0mm, % Sharp corners
    boxsep=1.0in, % space between box edges and text
    sidebyside, % Divide left/right
    halign=center,
    valign=center,
}
...
\begin{document}
...
\begin{tcolorbox}[<options>]
    ...
    \tcblower % divide box into two sections, upper and lower
    % or left and right if sidebyside option is set.
    ...
    \tcbsubtitle[<options>]{My subtitle}
\end{tcolorbox}
\end{lstlisting}

Example:
\begin{lstlisting}
\usepackage{tcolorbox}
\begin{tcolorbox}[colback=red!5!white,colframe=red!75!black,title=My nice heading]
    My awesome color box.
\end{tcolorbox}
\end{lstlisting}

\clearpage
\section{Minipage/columns}
\begin{lstlisting}
\begin{minipage}[<vertical align>][<height>]{<width>}
\raisebox{0pt}[<height>][<depth>]{...}
\begin{minipage}[t]{0.2\textwidth}
    stuff
\end{minipage}
\begin{minipage}[t]{0.8\textwidth}
    longer stuff
\end{minipage}
\end{lstlisting}

Use \verb|multicol| package:
\begin{lstlisting}
\usepackage{multicol}

% Start 2-columns
\begin{multicols}{2}
    ...
\end{multicols}

% Add asterisk to prevent forcing columns to equal heights
\begin{multicols*}{2}
    ...
\end{multicols*}

% No forcing cols to fill vertical space (same as above?)
\raggedcolumns

[
    \section{First section}
    Text that is not confined to declared columns. Not sure why you wouldn't
    just put this before starting the columns, but whatev.
]

\vfill                  % No forcing cols to fill vertical space (not working)
\columnbreak            % Start at top of next column

\addtolength{\columnsep}{5mm}  %  add space between columns
\setlength{\columnseprule}{0.4pt}  % set thickness of line between columns

% 02 October 2018:
\setlength{\columnsep}{0.20in}
    % Default = 28.6 pt. Divide by 72 to get inches.

\end{lstlisting}

\clearpage
\section{Symbols}
\begin{lstlisting}
\AA{}   % Angstrom (does not go between $s)
\infty  % infinity
\sim    % '~'
\approx % 'double ~'
\propto % proportionality symbol (like alpha)
\equiv  % like '=', but with three lines.
\& \%   % include these symbols in document
        % (also precede a space with '\' when in math mode).
\pm     % plus or minus (\mp for minus or plus)
\textbackslash  % \
\textgreater    % >
\textless       % <
\end{lstlisting}

$\textrm{some\ text}$    %   Roman style (removes italics),
                         %   and '\' puts space between words

\section{Lines}
\begin{lstlisting}

\hline % forces a break between paragraphs
\rule{length}{thickness} % Doesn't force break between paragraphs
\line(x-slope, y-slope){length}  % Syntax

% Example of horizontal line.
% Can't find an explanation of units for 'length'.
\line(1,0){450}
\dotfill
\hrulefill
\end{lstlisting}


\clearpage
\section{verbatim}
\begin{lstlisting}
\verb|text|  % inline

\begin{verbatim}  % environment
    ... text ...
\end{verbatim}
\end{lstlisting}

\section{lstlisting}
\url{http://users.ecs.soton.ac.uk/srg/softwaretools/document/start/listings.pdf}

Use to write code into latex documents (nicer than verbatim).
\begin{verbatim}
\usepackage{listings}

\lstset{ %
  backgroundcolor=\color{white},   % need package color or xcolor
  basicstyle=\footnotesize,        % font size
  breakatwhitespace=false,         % automatic breaks only at whitespace
  breaklines=true,                 % automatic line breaking
  captionpos=b,                    % sets the caption-position to bottom
  commentstyle=\color{mygreen},    % comment style
  deletekeywords={...},            % delete keywords from the given language
  escapeinside={\%*}{*)},          % for adding LaTeX within your code
  extendedchars=true,              % non-ASCII characters; 8-bits encodings only;
                                   %   does not work with UTF-8
  frame=single,                    % adds a frame around the code
  keepspaces=true,                 % keeps spaces in text (e.g. indentation)
                                   %  (possibly needs columns=flexible)
  keywordstyle=\color{blue},       % keyword style
  language=Octave,                 % the language of the code
  otherkeywords={*,...},           % if you want to add more keywords to the set
  numbers=left,                    % where to put the line-numbers;
                                   %   possible values are (none, left, right)
  numbersep=5pt,                   % how far the line-numbers are from the code
  numberstyle=\tiny\color{mygray}, % the style that is used for the line-numbers
  rulecolor=\color{black},         % if not set, the frame-color may be changed
                                   %  on line-breaks within not-black text
                                   %  (e.g. comments (green here))
  showspaces=false,                % show spaces everywhere adding particular
                                   %   underscores;
                                   %   it overrides 'showstringspaces'
  showstringspaces=false,          % underline spaces within strings only
  showtabs=false,                  % show tabs within strings adding
                                   %   particular underscores;
  stepnumber=2,                    % the step between two line-numbers.
                                   %  If it's 1, each line will be numbered
  stringstyle=\color{mymauve},     % string literal style
  tabsize=2,                       % sets default tabsize to 2 spaces
  title=\lstname                   % show the filename of files included
                                   %  with \lstinputlisting;
                                   %  also try caption instead of title
}
...
\begin{lstlisting}
    code code code
\end{lstlisting}
\end{verbatim}


\section{Environments and commands}

\textbf{Defining a new command}

Syntax: \verb|\newcommand{<cmd>}[<n>][<opt>]{<stuff>}|
\begin{description}
    %\item [\verb|n|] Number of arguments
    %\item [\verb|opt|] Options
    %\item [\verb|stuff|] stuff
    \item [n] Number of arguments
    \item [opt] Options
    \item [stuff] stuff
\end{description}
Existing environments (list, adjustwidth, etc.) can be used inside new commands.

\textbf{Changing an existing environment}
\begin{lstlisting}
\renewenvironment{name}{%
    ...}
\end{lstlisting}

\textbf{Defining a new environment}
\begin{lstlisting}
\newenvironment{name}[#]{%
    {<initialization code> (before text)}
    {<finalization code> (after text)}
}
\end{lstlisting}

\clearpage
\section{Graphics}
\subsection{Figures}
\begin{lstlisting}
\usepackage{graphicx}
  % Not needed with Beamer?
  % Seems to be needed for \graphicspath:
\graphicspath{ % If only using one path, still need both sets of {}s
    {/first/path/to/graphics/}
    {/second/path/to/graphics/}
    ...
}
\usepackage{float} % manage floating graphics
...
\begin{figure}[<placement specifier(s)>]
\centering
\includegraphics[<options>]{GreekSymbols.jpg}
\caption{How to insert greek symbols in LaTeX}
\label{greek}
\end{figure}
\end{lstlisting}
Options:
\begin{itemize}
    \item width=5.0in
    \item draft=true|false
    \item angle=90|180|etc\ldots  % rotate graphic counter-clockwise
\end{itemize}

Note: use a tilde, e.g.
\verb|See figure~\ref{figlabel}|
to prevent ``figure'' and number from being separated at a line break.

Placement specifiers:
\begin{description}
    \item [t] Top
    \item [b] Bottom
    \item [p] Page of floats
    \item [h] Here, if possible
    \item [H] Here, definitely
    \item [!] Force LaTeX to ignore its own determined location
\end{description}
LaTeX thinks it knows where to put your figures better than you do\ldots

Note that placement specifiers are for floating figures. If you're using, e.g.
the deluxetable environment with aastex, these options won't work (you'll
actually get an error).

\subsection{Tables}
\begin{verbatim}
\renewcommand{\arraystretch}{2}  % Apply doubling spacing between table rows, if desired
\setlength{\tabcolsep}{12pt}  % Add space between columns
\begin{table}[h]
    \caption{Values for polytropic index $n$ = 4.5}
    \centering
    \begin{tabular}{ c c c c c c c c c  }
        \hline\hline
        $n$ & $\xi_1$ & $\rho_c/\rho$ & $N_{n}$ & $W_n$ & $\Theta_n$
        & $\rho_c[g\,cm^{-3}]$ & $P_c[dyne\,cm^{-2}]$ & $T_c[K]$ \\
        \hline
        4.5 & 31.841 & 6187.500 & 0.658 & 4917.415 & 3.329 & 8718.704 &
        5.535e19 & 4.742e7 \\
        \hline
    \end{tabular}\\
    \label{table:nonlin}
\end{table}
\end{verbatim}
For the \verb|tabular| line, \verb|c| stands for center-justified;
use \verb|l| and \verb|r| for left and right justified.


\begin{lstlisting}
\begin{tabular{r p{6in}}
    one & two \newline more text
\end{tabular}
  % p = top alignment
  % b = bottom alignment
  % m = center alignment
\end{lstlisting}
The \verb|p/b/m| options are ``paragraph options''. They're left-aligned, let
you set the width of the cell so that long text will wrap nicely, and allow the
use of \verb|\newline| in the tabular environment, if needed.
\begin{minipage}{0.6\textwidth}
\begin{itemize}
    \item \verb|A \rotatebox{90}{B}C| \hfill A \rotatebox{90}{B}C
    \item \verb|A \rotatebox{270}{B}C| \hfill A \rotatebox{270}{B}C
    \item \verb|A \rotatebox[origin=c]{270}{B}C| \hfill A \rotatebox[origin=c]{270}{B}C
\end{itemize}
\end{minipage}

\subsection{Background}

Paragraph sep is \the\parskip

\begin{lstlisting}
\usepackage[<options>]{background}
 % or
\usepackage{background}
\backgroundsetup{<options>}  % Can be used in body of document

% Options:
pages=all|some
opacity=n  % 0 <= n <= 1
color=
contents=

\BgThispage
\NoBgThispage
\end{lstlisting}


\subsection{Grids and diagrams}
\begin{lstlisting}
\usepackage[<options>]{tikz}
...
\begin{document}

% Dimensions to use for entire pic
% [cm...where tikzpicture options define the units?]
\def\w{18}
\def\h{18}

\begin{tikzpicture}[x=1cm, y=1cm, semitransparent]
%  x,y -> units for grid coordinates below

    \draw[step=1mm, line width=0.1mm, black!30!white] (0,0) grid (\w, \h);
    % step - width of each box
    % line width - thickness of lines
    % (0,0) - relative to... ?

    \node[draw] at (0,0) {text}  % align = left|right, text width=

\end{lstlisting}



\section{Bibliographies}
Bibtex - entries are stored in a separate file, e.g. \verb|reffile.bib|, then
imported into the main document.
This file is formatted like \verb|@article{id,...}|.
(NOTE: Bibtex is NOT a package. Also not specific to aastex.)
\begin{lstlisting}
\bibliographystyle{plain}
\bibliographystyle{plainnat}  % Comes with natbib package; displays full url
\bibliographystyle{te} % te - one of many formatting styles; optional?
\begin{document}
... \cite{id} ...
\bibliography{reffile}
\end{document}
\end{lstlisting}

Natbib package:
\begin{lstlisting}
\usepackage{natbib}
%% In text citations:
\cite{Smith2018} % --> Smith et al. (2018)
\citep{Smith2018} % --> (Smith et al. 2018)
\cite{label1, label2} % 1+ papers by same author
\citet[p.~199]{label} % cite specific page
\citealt{label} % ?

\end{lstlisting}

\section{Labels and cross-referencing}
\begin{lstlisting}
\label{sec:labelname}...\ref{ssec:labelname}   % section
\label{ssec:labelname}...\ref{ssec:labelname}  % subsection
\label{sssec:labelname}...\ref{ssec:labelname} % subsubsection
\label{fig:labelname}...\ref{fig:labelname}    % Figures
\label{tab:labelname}...\ref{tab:labelname}    % Tables
\end{lstlisting}
In case same name is used for multiple things.
Also requires multiple runs of pdflatex.

\clearpage
\section{Maths!}
\begin{itemize}
    \item \href{http://www.math.harvard.edu/texman/node17.html}{Link}
    \item \href{http://www.math.illinois.edu/~ajh/tex/displays.html}{Link}
\end{itemize}

\subsection{In-text math mode}
Examples
\begin{itemize}
    \item \verb|$\frac{1}{4}$| $\rightarrow$ $\frac{1}{4}$
    \item \verb|$G=6.67\times10^{-8}$| $\rightarrow$
        $G=6.67\times10^{-8}$
\end{itemize}
If text is bold, make math symbols bold as well:
\begin{verbatim}
\usepackage{amsmath}
...
\textbf{This article discusses the $\beta$ parameter}
\textbf{This article discusses the $\boldsymbol\beta$ parameter}
\end{verbatim}
\textbf{This article discusses the $\beta$ parameter}\\
\textbf{This article discusses the $\boldsymbol\beta$ parameter}

\subsection{Equations}
\subsubsection{Numbered equations}
\begin{verbatim}
\begin{equation}
    P_{\textrm{mag}} = \frac{B^2}{\sqrt{4\pi\rho_o}}
\end{equation}
\end{verbatim}

\begin{equation}
    \boxed{P_{\textrm{mag}} = \frac{B^2}{\sqrt{4\pi\rho_o}}}
\end{equation}



\subsubsection{Equations without numbering}
Note that the \verb|\boxed{...}| commands are putting the examples
in boxes, but are not necessary for writing equations.
\begin{verbatim}
\begin{equation*}
    \boxed{%
    P_{\textrm{mag}} = \frac{B^2}{\sqrt{4\pi\rho_o}}
    }
\end{equation*}
\end{verbatim}
\begin{equation*}
    \boxed{%
    P_{\textrm{mag}} = \frac{B^2}{\sqrt{4\pi\rho_o}}
    }
\end{equation*}

Or simply put double \verb|$|s on each side of equation:
\begin{lstlisting}
$${
    P_{\textrm{mag}} = \frac{B^2}{\sqrt{4\pi\rho_o}}
}$$
\end{lstlisting}
$$  P_{\textrm{mag}} = \frac{B^2}{\sqrt{4\pi\rho_o}} $$
This may not work for more complicated math, such as matrices.
Apparently it is now best to use brackets rather than \$\$s:
\begin{lstlisting}
\[
    P_{\textrm{mag}} = \frac{B^2}{\sqrt{4\pi\rho_o}}
\]
\end{lstlisting}

\subsubsection{Aligning equations}
\begin{verbatim}
\usepackage{amsmath}
...
\begin{align}
k_1 &= hf(x_n,y_n)\\
k_2 &= hf(x_n+\frac{1}{2}h,y_n+\frac{1}{2}k_1)\\
k_3 &= hf(x_n+\frac{1}{2}h,y_n+\frac{1}{2}k_2)\\
k_4 &= hf(x_n+h,y_n+k_3)\\
y_{n+1} &=
y_n+\frac{1}{6}k_1+\frac{1}{3}k_2+\frac{1}{3}k_3+\frac{1}{6}k_4+O(h^5)\\
\end{align}
\end{verbatim}
\begin{align}
k_1 &= hf(x_n,y_n)\\
k_2 &= hf(x_n+\frac{1}{2}h,y_n+\frac{1}{2}k_1)\\
k_3 &= hf(x_n+\frac{1}{2}h,y_n+\frac{1}{2}k_2)\\
k_4 &= hf(x_n+h,y_n+k_3)\\
y_{n+1} &=
y_n+\frac{1}{6}k_1+\frac{1}{3}k_2+\frac{1}{3}k_3+\frac{1}{6}k_4+O(h^5)\\
\end{align}

Can also remove numbering from aligned equations:
\begin{verbatim}
\begin{align*}
    ...
\end{align*}
\end{verbatim}

\subsection{Size of brackets, parentheses, etc.}
In order of increasing size:
\begin{verbatim}
\big( ... \big)
\Big( ... \Big)
\bigg( ... \bigg)
\Bigg( ... \Bigg)
\end{verbatim}
BETTER\@:
\begin{verbatim}
\left( ... \right)
\end{verbatim}
to scale size of brackets to what is inside them!

Increase size of fraction inside text:
\begin{verbatim}
\cfrac{1}{2}
\end{verbatim}

There are $\frac{1}{2}$ as many as there were.\par
There are $\cfrac{1}{2}$ as many as there were.

\subsection{Arrays}
\begin{lstlisting}
\[ \left\{
    \begin{array}{l c r c l}
        x & y & z \\
    \end{array}\right.
    \]
\end{lstlisting}

\subsection{Superscripts, subscripts, and prescripts}
\begin{lstlisting}
\[
    \sum_{\mathclap{j=1}}x  % \usepackage{mathtools}
    \prescript{238}{92}{U}
    \]
\end{lstlisting}


\[
    \sum_{\mathclap{j=1}}x \qquad
    \prescript{238}{92}{U}
    \]

\subsection{Referring to parts of equation}
\url{http://tex.stackexchange.com/questions/261315/how-to-change-color-of-underbrace}
\begin{lstlisting}
\usepackage{amsmath}
...
\begin{document}
...
\<command>[<width>][<depth>]{<stuff>}
\end{lstlisting}
Possible commands:
\begin{itemize}[label={}]
    \item underbrace
    \item overbrace
    \item underbracket
    \item overbracket
\end{itemize}

\begin{lstlisting}
\usepackage{mathtools}
\usepackage{ragged2e}
\newlength\ubwidth
\newcommand\parunderbrace[2]{%
    \settowidth\ubwidth{$#1$}
    \underbrace{#1}_{\parbox{\ubwidth}{\scriptsize\RaggedRight#2}}}
\end{lstlisting}

Example:
\begin{lstlisting}
$\underbrace{P(X \mid O)}_{p_1} \propto \overbrace{P(X)P(O \mid X)}^{p_2}$\\
$\underbrace{P(X \mid O)}_{\text{This explains this part}}
\propto \overbrace{P(X)P(O \mid X)}^{\text{And this explains the other part}}$
\end{lstlisting}

%\begin{tcolorbox}  % Desktop doesn't like this
\begin{framed}
$\underbrace{P(X \mid O)}_{p_1} \propto \overbrace{P(X)P(O \mid X)}^{p_2}$\\
$\underbrace{P(X \mid O)}_{\text{This explains this part}}
\propto \overbrace{P(X)P(O \mid X)}^{\text{And this explains the other part}}$
\end{framed}
%\end{tcolorbox}  % Desktop doesn't like this


\subsection{Operations}
\subsubsection{Integrals}
\begin{verbatim}
$\int$ % indefinite integral
$\int_{x1}^{x2}$ % definite integral, between x1 and x2
\end{verbatim}
\subsubsection{Square root}
\begin{verbatim}
$\sqrt{2\ln(2)}$
\end{verbatim}
\subsubsection{Summation (and the multiplication version)}
\begin{verbatim}
$$\sum_{n=1}^{\infty} 2^{-n} = 1$$
\end{verbatim}
$$\sum_{n=1}^{\infty} 2^{-n} = 1$$

\begin{verbatim}
$$ P(D|M) \propto \prod^{N-1}_{i=0}\left\{\exp
    \left[-\frac{1}{2}\left[\frac{y_i-y(x_i|a_j)}
{\sigma}\right]^2\right]\Delta y \right\} $$
\end{verbatim}
$$ P\left(D|M\right) \propto \prod^{N-1}_{i=0}\left\{\exp
    \left[-\frac{1}{2}\left[\frac{y_i-y\left(x_i|a_j\right)}
{\sigma}\right]^2\right]\Delta y \right\} $$


\clearpage
\section{Including separate files}
Files in current directory: \texttt{file1.tex, file2.tex, file3.tex}.
Then edit main file (e.g. \texttt{main.tex}):
\begin{lstlisting}
% To only include one or a few files (everything in main file is compiled).
\includeonly{file1, file3}
\begin{document}
...
\include{file1}
...
\include{file2}
...
\include{file3}
...
\end{lstlisting}


\begin{itemize}
    \item \verb|input|
        \begin{itemize}
            \item No extra files created
            \item No `input only' option
            \item No page breaks around the file contents
            \item Can be used inside files to be input/included.
        \end{itemize}
    \item \verb|include|
        \begin{itemize}
            \item Creates second file \texttt{file.aux}, which keeps track
                of counter values to preserve page numbers, chapters, etc.
            \item Can use `includeonly' to keep from compiling everything
            \item Effectively places \verb|\clearpage| before and after the
                file contents
            \item Can't be used inside files to be input/included.
        \end{itemize}
\end{itemize}

\verb|\include{file.tex}| does nothing.
\verb|\input{file.tex}| works, but may as well be consistent, and never
include file extensions here either.




\clearpage %--------------------------------------------------------------------
\section{Unsorted}


\begin{lstlisting}
% Print the numerical value of a latex length (in body text):
\usepackage{printlen}
% \begin{document}
% ...
% \printlength{length}
%   units specified by \uselengthunit{unit},
%   'unit' = pt, pc, in, mm, cm, bp, dd, cc
\end{lstlisting}



Define custom width in preamble if used many times in document, and you
don't want to manually change every single line if you decide to adjust
the width.
\begin{lstlisting}
\newlength{\mylength}
\setlength{\mylength}{0.8\textwidth}
\end{lstlisting}



Check \texttt{filename.log} for version of packages used. May need to add
\verb|\listfiles| in the preamble first.

\begin{lstlisting}
\usepackage{lipsum}
\lipsum[2-4] % paragraphs 2 -> 4
\end{lstlisting}
\begin{enumerate}
    \item shortest
    \item second-longest
    \item second-shortest
    \item Longest
\end{enumerate}

Counters:
\begin{lstlisting}
\newcounter{paranum}
\refstepcounter{paranum}\theparanum
 % Before adding the 'refstepcounter{paranum} part, the counter started at 0
 %  and didn't increment. Adding it started at 1, and did increment.

\the<counter>  % In general...

% Levels to show in table of contents
\setcounter{tocdepth}{4}

% Levels of sections to show numbers (4 --> paragraphs are numbered)
\setcounter{secnumdepth}{4}
\end{lstlisting}

\begin{lstlisting}
    \usepackage{amsmath, amsfonts}
    $ a\:\text{and}\:b_\in_\mathbb{N}$

\end{lstlisting}
\[
    a\:\text{and}\:b_{\in_{\mathbb{N}}}
    \]


In \texttt{think\_python.tex}, add part for using straight single quotes
in verbatim environment.


Lined background (like notebook paper):
\begin{lstlisting}

    \documentclass{article}
    \usepackage[vmargin=3cm]{geometry}
    \usepackage{tikzpagenodes}
    \usepackage{lipsum}
    \usepackage{background}

    \definecolor{notepadrule}{RGB}{217,244,244}

    \backgroundsetup{
    contents={%
      \begin{tikzpicture}
        \foreach \fila in {0,...,52}
        {
          \draw [line width=1pt,color=notepadrule]
          (current page.west|-0,-\fila*12pt) -- ++(\paperwidth,0);
        }
        \draw[overlay,red!70!black,line width=1pt]
          ([xshift=-1pt]current page text area.west|-current page.north) --
          ([xshift=-1pt]current page text area.west|-current page.south);
      \end{tikzpicture}%
    },
    scale=1,
    angle=0,
    opacity=1
    }

    \begin{document}

    \lipsum[1-14]


\end{lstlisting}


Grid:

\begin{lstlisting}
    \documentclass{article}
    \usepackage{tikz}
    \usepackage{multido}  % ?
    \usepackage[margin=0pt]{geometry}  % To cover entire page
    \unitlength=1mm

    \definecolor{purple}{rgb}{0.44, 0.16, 0.39}
    \definecolor{green}{rgb}{0.12, 0.3, 0.17}
    \definecolor{blue}{rgb}{0.03, 0.27, 0.49}
    \definecolor{cyan}{rgb}{0.0, 0.55, 0.55}


    \begin{document}


    % Graph paper
    %\newgeometry{left=1.0pt}
    \begin{tikzpicture}[x=1.0in, y=1.0in] %, semitransparent]
        \draw[step=0.25in,line width=0.2mm,blue!20!white]
        (0,0) grid (8.5,11.0);
    \end{tikzpicture}
    \newpage
    \begin{tikzpicture}[x=1.0in, y=1.0in] %, semitransparent]
        \draw[step=0.25in,line width=0.2mm,blue!20!white]
        (0,0) grid (8.5,11.0);
    \end{tikzpicture}


    % Dotted grid
    \newpage
    \begin{tikzpicture}[scale=.5]
        \foreach \x in {0,...,43}
        \foreach \y in {0,...,55}{\fill[black!20!white] (\x,\y) circle (0.06cm);}
      \end{tikzpicture}
    \newpage
    \begin{tikzpicture}[scale=.5]
        \foreach \x in {0,...,43}
        \foreach \y in {0,...,55}{\fill[black!20!white] (\x,\y) circle (0.06cm);}
      \end{tikzpicture}

    %\restoregeometry
\end{lstlisting}


\newpage

\begin{lstlisting}
%% outline-sample.tex
%% Copyright 1991 Peter Halvorson
%% Updates for LaTeX2e copyright 2002 Seth Flaxman
%% Updated for LPPL 1.3c or later by Clea F. Rees (for Seth Flaxman), 2008/10/06.
%
% This work may be distributed and/or modified under the
% conditions of the LaTeX Project Public License, either version 1.3
% of this license or (at your option) any later version.
% The latest version of this license is in
%   http://www.latex-project.org/lppl.txt
% and version 1.3 or later is part of all distributions of LaTeX
% version 2005/12/01 or later.
%
% This work has the LPPL maintenance status `unmaintained'.
%
% This work consists of the files outline.sty and outline-sample.tex.
% Save file as: outline-sample.tex

\documentclass{report}
\usepackage{outline}

% [outline] includes new outline environment. I. A. 1. a. (1) (a)
% use \begin{outline} \item ... \end{outline}

\pagestyle{empty}

\begin{document}

\begin{outline}
  \item {\bf Introduction }
  \begin{outline}
    \item {\bf Applications } \\
      Motivation for research and applications related to the
      subject.
    \item {\bf Organization } \\
      Explain organization of the report, what is included, and what
      is not.
  \end{outline}
  \item {\bf Literature Survey }
  \begin{outline}
    \item {\bf Experimental Work } \\
      Literature describing experiments with something in common with
      my experiment.  My experiment is subdivided into section
      relating to each aspect of the whole.
    \begin{outline}
      \item {\bf Drop Delivery } \\
    Literature relating to the production of droplets.
      \begin{outline}
    \item {\bf Continuous } \\
      Continuous drop production methods, i.e. jet methods.
    \item {\bf Drop on Demand } \\
      Drop on demand methods, i.e. ink jet devices.  Produce drops
      whenever needed, simplifies control of frequency.
    \item {\bf Flexibility } \\
      Best methods in terms of flexible velocities, volumes, and
      frequencies.
    \item {\bf Control Circuitry } \\
      Circuitry necessary to control the drops, may include
      control of generation, size, and frequency.  Divertors and
      drop chargers.
    \item {\bf Extensibility } \\
      Methods extensible to 2D applications.
    \item {\bf Recirculation } \\
      Recirculation techniques, pump, none, capillary.
      \end{outline}
      \item {\bf Instrumentation } \\
    Literature dealing with measurement of various parameters.
      \begin{outline}
    \item {\bf Temperature }
        \begin{outline}
          \item {\bf Heater Surface }
      \item {\bf Fluid Temperature }
      \item {\bf Heat Flux }
      \item {\bf Heat Transfer Coefficient }
        \end{outline}
    \item {\bf Drop Characteristics }
    \begin{outline}
      \item {\bf Size }
      \item {\bf Velocity }
      \item {\bf Frequency }
        \end{outline}
      \end{outline}
      \item {\bf Heating Element } \\
    Literature dealing with the heating element.  Material
    properties, surface properties, heat sources.
      \begin{outline}
    \item {\bf Material }
    \item {\bf Heat Source }
      \end{outline}
    \end{outline}
    \item {\bf Analytical Work }
    \begin{outline}
      \item {\bf Evaporation }
      \item {\bf Boiling }
      \item {\bf Leidenfrost Temperatures }
      \item {\bf Heat Transfer }
      \item {\bf Numerical Analysis }
      \begin{outline}
    \item {\bf Drop Characteristics }
    \item {\bf Surface Wetting }
    \item {\bf Transient Temperatures }
      \end{outline}
    \end{outline}
  \end{outline}
  \item {\bf Proposed Research }
  \begin{outline}
    \item {\bf Experimental Work }
    \item {\bf Analytical Work }
  \end{outline}
\end{outline}

\end{document}
\end{lstlisting}



\end{document}
